\section{Introduction}

Recent advances in generative AI have changed the organizational problem faced by firms. The relevant technology is not merely an internal productivity tool. Increasingly, firms access powerful AI systems through external platforms and use them for multi-step execution rather than one-shot prediction. As a result, the economic issue is not simply whether some tasks can be automated, but how reliance on externally supplied agentic AI reshapes the internal allocation of authority and the preservation of organizational capabilities \citep{AgrawalGansGoldfarb2019,NzembayieUrbano2026}.

In such an environment, the firm must decide which tasks to assign to external AI systems and which tasks to retain under human authority. This question belongs to the classic economics of organizations, where authority, communication, and the distribution of knowledge shape internal decision-making \citep{AghionTirole1997,Dessein2002,Garicano2000}. It is also related to recent work on the allocation of decision authority between humans and artificial intelligence \citep{AtheyBryanGans2020}. However, the organizational problem created by current generative AI differs from earlier formulations in one important respect: the firm delegates not to an internal machine whose use is fully under its control, but to an external platform whose services are priced and whose availability cannot be taken for granted.

This paper develops a simple two-period model of authority allocation with external agentic AI. In period 1, the firm assigns routine tasks to an external AI platform and retains sufficiently exceptional tasks under human authority. In period 2, contingencies may arise in which the external platform becomes unavailable or cannot be fully relied upon, so that the firm must draw on retained internal recovery capability (fallback capability). The model therefore links current authority allocation to future organizational resilience. Relative to the existing literature, I treat AI not as a one-shot signal or decision aid, but as an externally supplied multi-step execution technology, and I allow current delegation choices to affect the firm's future capacity to recover, override, and coordinate when external AI fails or becomes unusable \citep{AgrawalGansGoldfarb2019,AtheyBryanGans2020,NzembayieUrbano2026}.

The analysis yields three main results. First, the optimal organization takes a hybrid form: routine tasks are delegated to AI, while sufficiently exceptional tasks remain under human authority. Second, full automation may be suboptimal even when AI dominates human managers task by task in normal operations. The reason is not that AI necessarily performs poorly on marginal tasks, but that extensive reliance on external AI erodes retained human recovery capability that becomes valuable under contingencies. Third, platform pricing distorts the firm's internal authority structure relative to a technological benchmark, and this distortion is stronger in more agentic environments with longer execution horizons. More broadly, the paper shows that a new externally supplied execution technology can alter not only current task allocation, but also the firm's internal design of control, adaptation, and resilience.

The paper contributes to three strands of literature. It contributes to the economics of authority, communication, and organizational knowledge \citep{AghionTirole1997,Dessein2002,Garicano2000}; to recent work on AI and the allocation of decision authority \citep{AtheyBryanGans2020,AgrawalGansGoldfarb2019}; and to research on technical change, governance, and institutional dependence in platform environments \citep{RochetTirole2003,BakerGibbonsMurphy2023,MarschakWei2024,MukherjeeSaha2024,Kim2025,NzembayieUrbano2026}. In particular, the paper speaks to a JITE-style question: how a new externally supplied technology interacts with internal governance design, adaptation, and organizational control. The central message is that when critical AI capabilities are supplied by external platforms, the firm faces a distinct organizational trade-off between current execution efficiency and the preservation of internal recovery capability.

The remainder of the paper is organized as follows. Section~\ref{sec:lit} discusses the related literature. Section~\ref{sec:model} presents the model. Sections~\ref{sec:baseline}--\ref{sec:capability} derive the main results. Section~\ref{sec:discussion} discusses broader institutional implications, and Section~\ref{sec:conclusion} concludes.
